% JAC1964-20261121. /root/programme_audit; GPT-6 Astra Pro.
% P17: actual rational Chow descent, not only realization descent.

\subsection*{JAC-P17. The Fermat motive and cycle over their precise fields}

\begin{proposition}[JAC-P17]\label{jac:P17}
Let $C$ be either of the two rational Fermat models
\[
 C_-:x^4+y^4=z^4,\qquad C_+:x^4+y^4+z^4=0,
\]
let $J$ be its Jacobian, and let $G=\mu_4^2$ act by the two
coordinate scalings relative to $z$. Put $F=\mathbf Q(i)$.
The exact Chow motive
\[
                    N=(h^3(J),p_G)
\]
of JAC-P13 is defined over $\mathbf Q$. The Fourier--Ceresa
class $\eta$ defined using a hyperflex is an actual class in
$\operatorname{CH}^2(N)_{\mathbf Q}$ over $\mathbf Q$.
The two models give canonically identified pairs $(N,\eta)$
over $\mathbf Q$ under their coordinate-scaling isomorphism.

There is an explicit $F$-defined central endomorphism
$j\in\operatorname{End}^0(J_F)$ with $j^2=-1$, whose normalized
middle pullback is the operation $T$ of JAC-P11--P14.
If $c$ is the nontrivial automorphism of $F/\mathbf Q$, then
\[
                  c(j)=-j,\qquad c(T)=-T.              \tag{JAC.P17.1}
\]
Write $V_L=\operatorname{CH}^2(N_L)_{\mathbf Q}$ for a field
$L\supset\mathbf Q$. Restriction gives the exact decomposition
\[
                  V_F=V_{\mathbf Q}\oplus T V_{\mathbf Q},
                                                               \tag{JAC.P17.2}
\]
where the summands are respectively the $+1$ and $-1$
eigenspaces of $c$. Equivalently, $V_F$ is the scalar extension
of the rational vector space $V_{\mathbf Q}$ to $F$, with
$i$ acting through $T$. This is an assertion about the whole
Chow group of $N$, not just the known Ceresa plane.

For the degree-normalized height used in AR-2 and JAC-P11,
the two summands in (JAC.P17.2) are orthogonal and
\[
 h_F(z+Tw,z+Tw)=h_{\mathbf Q}(z,z)+h_{\mathbf Q}(w,w)
       \quad(z,w\in V_{\mathbf Q}).                   \tag{JAC.P17.3}
\]
In particular positivity, semipositivity, or zero radical on
the entire $V_F$ is equivalent to the respective assertion on
$V_{\mathbf Q}$. The known positive plane is
$\mathbf Q\eta\oplus\mathbf Q T\eta$; its rational invariant
part is exactly $\mathbf Q\eta$.
No assertion that this line equals $V_{\mathbf Q}$ is made.
\end{proposition}

\begin{proof}
\emph{The projector and the cycle.}
The Deninger--Murre degree-three projector $\pi_3$ of $J$ is
defined over $\mathbf Q$ and is functorial for its endomorphisms.
Over $F$, form the actual Chow correspondence
\[
          p_N=\pi_3\,\frac1{16}\sum_{g\in G(\overline{\mathbf Q})}g_*.
                                                               \tag{JAC.P17.4}
\]
Conjugation permutes the sixteen terms, replacing the two
coordinate generators by their inverses. Thus the sum is the
base change of a rational correspondence. One may construct
it explicitly as half the finite pushforward from $F$ of the
displayed correspondence. Its pullback is (JAC.P17.4).
The relations $p_N^2=p_N$ and $g_*p_N=p_Ng_*=p_N$ hold
exactly over $F$ by averaging, and descend to $\mathbf Q$:
restriction on rational Chow groups under a finite field
extension is injective because norm composed with restriction
is multiplication by its degree. Consequently this is an actual
Chow-motive descent. JAC-P13 identifies this projector with the
complementary primitive motive used in the height calculation.

The principal polarization of a Jacobian is defined over its
ground field. Its symmetric rational divisor class $\theta$ is
therefore defined over $\mathbf Q$, as is the Fourier
correspondence attached to it. One way to see the divisor-class
assertion without assuming a rational theta divisor is to take
a symmetric representative over a finite extension and average
its conjugates. Two symmetric representatives of the same
polarization differ by symmetric classes in $\operatorname{Pic}^0$,
which are two-torsion and hence zero with rational coefficients.
Thus the averaging preserves its geometric class exactly.

Every hyperflex $P$ satisfies $4P\sim K_C$. For any Galois
conjugate $\sigma P$, the difference $\sigma P-P$ is therefore
four-torsion in $J$. As checked in AR-8, translation by a torsion
point is the identity on rational Chow groups: if $na=0$, then
$[n]_*t_{a*}=[n]_*$, and $[n]_*$ is invertible on rational Chow
groups by the abelian Chow--Kunneth decomposition.
It follows that the class of the Abel curve
$x\mapsto[x-P]$ is independent of the hyperflex and invariant
under Galois. All these equalities already hold over a finite
Galois field containing the hyperflexes in question.
Normalized norm from that field thus gives a class over
$\mathbf Q$ whose restriction is this very Abel-curve class.
Its grade-one projector and its Fourier transform are defined
over $\mathbf Q$. This gives the asserted actual class $\eta$.
The equality $p_N\eta=\eta$, already proved geometrically in
AR-8 and JAC-P13, descends by the same injectivity argument.
For $C_-$ one may simply use the rational hyperflex $(1:0:1)$;
for $C_+$ the hyperflex $[1:\zeta_8:0]$ gives the same
descent argument, with no rational-point assumption.

\emph{The explicit central operation.}
Let $u,v$ be the induced Jacobian automorphisms of
$(x:y:z)\mapsto(ix:y:z)$ and $(x:y:z)\mapsto(x:iy:z)$.
They are defined over $F$, commute, and preserve the canonical
principal polarization. Define, in the algebra of abelian
endomorphisms,
\[
 j=\frac14\bigl(u-u^{-1}+v-v^{-1}-uv+(uv)^{-1}\bigr).
                                                               \tag{JAC.P17.5}
\]
The three holomorphic characters are $(1,1),(1,2),(2,1)$.
On the first line, the numerator in (JAC.P17.5) has eigenvalue
$2i+2i-0=4i$; on the second and third lines it has eigenvalue
$2i+0-(-2i)=4i$. It has eigenvalue $-4i$ on the conjugate
lines. Therefore $j$ acts by $i$ on the holomorphic space and
by $-i$ on its conjugate.

A homomorphism of complex abelian varieties is determined by
its map on the tangent space, since its analytic lift is complex
linear. This proves $j^2=-1$ and centrality in the entire
geometric endomorphism algebra, not just in $\mathbf Q[u,v]$.
Because $u^\dagger=u^{-1}$ and $v^\dagger=v^{-1}$, the same
formula also gives $j^\dagger=-j$. Thus (JAC.P17.5) is exactly
the choice of central $j$ with holomorphic eigenvalue $i$ used
in JAC-P14.

Set $a=4j$, an actual $F$-endomorphism, and define
\[
                  T=4^{-3}a^*\bigm|_N.                \tag{JAC.P17.6}
\]
Functoriality and the homogeneous degree-three power formula
for abelian Chow motives, checked in JAC-P14, make this the
normalized pullback $\rho_3(j)$ independently of the denominator.
It commutes with $p_N$, has $T^2=-1$, and is the height isometry
of JAC-P11. Conjugation sends $u,v$ to $u^{-1},v^{-1}$ in
(JAC.P17.5), so $c(j)=-j$ exactly as an endomorphism.
The homogeneous odd-degree formula now gives
$c(T)=\rho_3(-j)=-\rho_3(j)=-T$ exactly in Chow motives.
No identification from realization alone is being used here.
Notice that $T$ is constructed by applying the degree-three
functor to a linear combination on $h^1$. A linear combination
of the original $G$-actions on $N$ itself would only be scalar,
because every element of $G$ acts as the identity there.

\emph{The two rational models.}
Over $L=\mathbf Q(\zeta_8)$, the map
\[
             (x:y:z)\longmapsto(x:y:\zeta_8z)
                                                               \tag{JAC.P17.7}
\]
is an isomorphism $C_-\to C_+$. For every
$\sigma\in\operatorname{Gal}(L/\mathbf Q)$, its conjugate
differs from it by scaling $z$ by
$\sigma(\zeta_8)/\zeta_8\in\mu_4$. Projectively this is an
element of the group generated by the $x$ and $y$ scalings.
The induced map between the $N$ summands is therefore invariant
under Galois as an actual Chow correspondence: the additional
factor is annihilated by the relation $gp_N=p_N$.
Its normalized norm descends the isomorphism to $\mathbf Q$;
the inverse descends in the same way, and their identities follow
by restriction injectivity. This isomorphism preserves the
canonical principal polarization and the hyperflex Abel-curve
class, and hence sends $\eta$ to $\eta$. It also commutes with
$u,v$ and therefore with (JAC.P17.5) and (JAC.P17.6).
Different choices of the fourth root of $-1$ differ by the same
trivial $G$ action on $N$.

\emph{Descent of the full Chow group and its height.}
For the quadratic extension $F/\mathbf Q$, restriction and
norm on the motive defined by (JAC.P17.4) satisfy
\[
 \operatorname{cor}\operatorname{res}=2,
 \qquad \operatorname{res}\operatorname{cor}=1+c.
                                                               \tag{JAC.P17.8}
\]
Thus restriction identifies $V_{\mathbf Q}$ with $V_F^{c=1}$.
The involution $c$ splits $V_F$ into its two eigenspaces over
$\mathbf Q$. Since $cT=-Tc$ and $T^{-1}=-T$, the map $T$
identifies the invariant space with the anti-invariant space.
This proves (JAC.P17.2), without a finite-dimensionality
assumption on these Chow groups.

The normalized height is invariant under conjugation and under
finite extension. Hence for invariant $z,w$,
\[
 h_F(z,Tw)=h_F(cz,cTw)=-h_F(z,Tw)=0.
\]
Together with the isometry property of $T$, this gives
(JAC.P17.3). It proves all three stated equivalences for the
whole groups. AR-13 and JAC-P11 imply that $\eta,T\eta$ are
linearly independent and have the equal positive norm; their
conjugation signs give the last assertion.
\end{proof}

\paragraph{The arithmetic character is unchanged by the model choice.}
The directly read BDP source, \emph{p-adic L-functions and the
coniveau filtration on Chow groups}, Section~5.2, printed
pp.~24--25, gives for $C_-$ the rational isogeny
$J\sim E_-^2\times E_+$, where $E_\pm:y^2=x^3\pm x$.
If $\psi_-$ and $\psi_+$ are their Gaussian Hecke characters,
then $\psi_+=\psi_-\mu$, with $\mu^2=1$ and conductor $(4)$.
The extreme weight-three character is
\[
                \chi=\psi_-^2\psi_+=\psi_+^3.
                                                               \tag{JAC.P17.9}
\]
This is the character used in AR-19. The descended isomorphism
of $N$ for the two models proves that this same induced
two-dimensional representation is the realization of the
$\mathbf Q$-motive in (JAC.P17.4), even when the stable-reduction
calculation used the plus model. This paragraph identifies the
representation and its fields; a bound for a Selmer group is
still not injectivity of the scalar-extended Chow Abel--Jacobi map.

